\section{Formulario}

\subsection{Definición de campo eléctrico}

El campo eléctrico se define como la fuerza eléctrica por unidad de carga de prueba positiva:

\[
\boxed{\vec E=\frac{\vec F_0}{q_0}}
\]

En magnitud,

\[
\boxed{E=\frac{F_0}{q_0}}.
\]

En componentes rectangulares,

\[
\boxed{E_x=\frac{F_x}{q_0}}, \qquad
\boxed{E_y=\frac{F_y}{q_0}}.
\]

Unidad SI:

\[
\boxed{[E]=\si{\newton\per\coulomb}}.
\]


\subsection{Fuerza sobre una carga en un campo eléctrico}

\[
\boxed{\vec F=q\vec E}
\]

En componentes,

\[
\boxed{F_x=qE_x}, \qquad
\boxed{F_y=qE_y}.
\]

Magnitud:

\[
\boxed{F=|q|E}
\]

Dirección:

\[
q>0 \Longrightarrow \vec F \text{ tiene el mismo sentido que }\vec E,
\]

\[
q<0 \Longrightarrow \vec F \text{ tiene sentido opuesto a }\vec E.
\]


\subsection{Campo eléctrico de una carga puntual}

Para una carga fuente puntual \(Q\), a una distancia \(r\),

\[
\boxed{E=k\frac{|Q|}{r^2}}
\]

con

\[
\boxed{k\approx\SI{8.99e9}{\newton\metre\squared\per\coulomb\squared}}.
\]

Dirección del campo:

\[
Q>0 \Longrightarrow \text{el campo apunta alejándose de la carga},
\]

\[
Q<0 \Longrightarrow \text{el campo apunta hacia la carga}.
\]

Dependencia con la distancia:

\[
\boxed{E\propto\frac{1}{r^2}}.
\]

Por tanto,

\[
r\to2r \Longrightarrow E\to\frac{E}{4},
\]

\[
r\to3r \Longrightarrow E\to\frac{E}{9},
\]

\[
r\to\frac{r}{2} \Longrightarrow E\to4E.
\]


\subsection{Campo de una carga puntual mediante coordenadas}

Si una carga fuente \(q_i\) está situada en

\[
(x_i,y_i)
\]

y queremos calcular el campo en

\[
P=(x,y),
\]

definimos

\[
\boxed{\Delta x_i=x-x_i}, \qquad
\boxed{\Delta y_i=y-y_i}.
\]

La distancia entre la carga y el punto es

\[
\boxed{r_i=\sqrt{(\Delta x_i)^2+(\Delta y_i)^2}}.
\]

La magnitud del campo es

\[
\boxed{E_i=k\frac{|q_i|}{r_i^2}}.
\]

Si se conoce el ángulo \(\theta_i\) del campo respecto del eje \(+x\),

\[
\boxed{E_{i,x}=E_i\cos\theta_i},
\qquad
\boxed{E_{i,y}=E_i\sin\theta_i}.
\]

Los signos de las componentes se determinan según el sentido real del campo en el plano.


\subsection{Componentes directas a partir de las coordenadas}

Las componentes también pueden calcularse directamente mediante

\[
\boxed{
E_{i,x}
=
kq_i\frac{x-x_i}{r_i^3}
}
\]

y

\[
\boxed{
E_{i,y}
=
kq_i\frac{y-y_i}{r_i^3}
},
\]

donde

\[
r_i=\sqrt{(x-x_i)^2+(y-y_i)^2}.
\]

Estas expresiones incorporan automáticamente el signo de la carga fuente.


\subsection{Principio de superposición}

El campo total es la suma vectorial de los campos producidos por todas las cargas.

Trabajando por componentes rectangulares,

\[
\boxed{E_x=\sum_i E_{i,x}},
\qquad
\boxed{E_y=\sum_i E_{i,y}}.
\]

La magnitud del campo resultante es

\[
\boxed{E=\sqrt{E_x^2+E_y^2}}.
\]

Para obtener un ángulo de referencia,

\[
\boxed{
\theta_r=
\tan^{-1}\left(\frac{|E_y|}{|E_x|}\right)
}.
\]

El cuadrante se determina con los signos de \(E_x\) y \(E_y\):

\[
\begin{array}{c|c|c}
E_x & E_y & \text{cuadrante}\\
\hline
+ & + & \text{I}\\
- & + & \text{II}\\
- & - & \text{III}\\
+ & - & \text{IV}
\end{array}
\]


\subsection{Superposición en una dimensión}

Si todos los campos están sobre el eje \(x\),

\[
\boxed{E_x=\sum_iE_{i,x}}.
\]

Si

\[
E_x>0,
\]

el campo resultante apunta hacia \(+x\).

Si

\[
E_x<0,
\]

apunta hacia \(-x\).


\subsection{Campo eléctrico nulo}

La condición general es

\[
\boxed{\vec E=\vec 0}.
\]

En componentes rectangulares,

\[
\boxed{E_x=0},
\qquad
\boxed{E_y=0}.
\]

En una dimensión,

\[
\boxed{\sum_iE_{i,x}=0}.
\]

Para dos contribuciones opuestas sobre una misma recta,

\[
\boxed{E_1=E_2}.
\]

Antes de igualar magnitudes debe comprobarse que los campos tengan sentidos opuestos.


\subsection{Dos cargas del mismo signo}

Sobre la línea que une dos cargas del mismo signo, el punto de campo nulo puede encontrarse entre ellas.

Si las cargas son \(q_1\) y \(q_2\), separadas una distancia \(L\), y \(x\) se mide desde \(q_1\),

\[
k\frac{|q_1|}{x^2}
=
k\frac{|q_2|}{(L-x)^2}.
\]

Por tanto,

\[
\boxed{
\frac{|q_1|}{x^2}
=
\frac{|q_2|}{(L-x)^2}
}.
\]

El punto se encuentra más cerca de la carga de menor magnitud.


\subsection{Dos cargas de signos opuestos}

Entre dos cargas de signos opuestos, los campos apuntan en el mismo sentido, por lo que

\[
\boxed{\text{no existe punto de campo nulo entre ellas}}.
\]

Si existe un punto de cancelación sobre la línea que las une, debe buscarse fuera del segmento que contiene las cargas.


\subsection{Distribuciones continuas de carga}

Para un elemento infinitesimal \(dq\),

\[
\boxed{dE=k\frac{|dq|}{R^2}}.
\]

El campo total se obtiene mediante integración:

\[
\boxed{\vec E=\int d\vec E}.
\]

Trabajando por componentes,

\[
\boxed{E_x=\int dE_x},
\qquad
\boxed{E_y=\int dE_y}.
\]


\subsection{Componentes de un elemento de campo}

Si el elemento \(dq\) produce un campo \(dE\) que forma un ángulo \(\theta\) con el eje \(+x\),

\[
\boxed{dE_x=dE\cos\theta},
\qquad
\boxed{dE_y=dE\sin\theta}.
\]

Los signos se asignan según el sentido de cada componente.


\subsection{Densidad lineal de carga}

\[
\boxed{\lambda=\frac{dq}{d\ell}}
\]

por lo que

\[
\boxed{dq=\lambda\,d\ell}.
\]

Unidad SI:

\[
\boxed{[\lambda]=\si{\coulomb\per\metre}}.
\]

Para una distribución uniforme de longitud \(L\),

\[
\boxed{\lambda=\frac{Q}{L}}.
\]


\subsection{Densidad superficial de carga}

\[
\boxed{\sigma=\frac{dq}{dA}}
\]

y

\[
\boxed{dq=\sigma\,dA}.
\]

Unidad SI:

\[
\boxed{[\sigma]=\si{\coulomb\per\metre\squared}}.
\]

Para una distribución uniforme sobre un área \(A\),

\[
\boxed{\sigma=\frac{Q}{A}}.
\]


\subsection{Densidad volumétrica de carga}

\[
\boxed{\rho=\frac{dq}{dV}}
\]

y

\[
\boxed{dq=\rho\,dV}.
\]

Unidad SI:

\[
\boxed{[\rho]=\si{\coulomb\per\metre\cubed}}.
\]

Para una distribución uniforme en un volumen \(V\),

\[
\boxed{\rho=\frac{Q}{V}}.
\]


\subsection{Campo de una distribución lineal}

Como

\[
dq=\lambda\,d\ell,
\]

la magnitud diferencial es

\[
\boxed{dE=k\frac{|\lambda\,d\ell|}{R^2}}.
\]

Por componentes,

\[
\boxed{E_x=\int dE_x},
\qquad
\boxed{E_y=\int dE_y}.
\]


\subsection{Campo de una distribución superficial}

Como

\[
dq=\sigma\,dA,
\]

\[
\boxed{dE=k\frac{|\sigma\,dA|}{R^2}}.
\]

El campo se obtiene integrando las componentes correspondientes.


\subsection{Campo de una distribución volumétrica}

Como

\[
dq=\rho\,dV,
\]

\[
\boxed{dE=k\frac{|\rho\,dV|}{R^2}}.
\]

El campo total se obtiene integrando sobre todo el volumen cargado.


\subsection{Campo sobre el eje de un anillo uniformemente cargado}

Para un anillo de radio \(a\), carga total \(Q\) y un punto situado a una distancia \(x\) sobre su eje,

\[
\boxed{
E_x=
k\frac{Qx}{(x^2+a^2)^{3/2}}
}.
\]

Por simetría,

\[
\boxed{E_y=0}.
\]

En el centro,

\[
x=0
\quad\Longrightarrow\quad
\boxed{E=0}.
\]

Si

\[
x\gg a,
\]

entonces

\[
\boxed{E\approx k\frac{|Q|}{x^2}},
\]

por lo que el anillo se comporta aproximadamente como una carga puntual vista desde muy lejos.


\subsection{Campo uniforme}

Un campo es uniforme cuando

\[
\boxed{\vec E=\text{constante}}
\]

en la región considerada.

Esto significa que su magnitud y dirección no cambian con la posición.

En componentes,

\[
\boxed{E_x=\text{constante}},
\qquad
\boxed{E_y=\text{constante}}.
\]


\subsection{Líneas de campo eléctrico}

Reglas principales:

\begin{itemize}
    \item Las líneas salen de cargas positivas.
    \item Las líneas entran en cargas negativas.
    \item La tangente a una línea indica la dirección de \(\vec E\).
    \item Una mayor densidad de líneas representa un campo más intenso.
    \item Dos líneas de campo no pueden cruzarse.
    \item En electrostática, las líneas de campo no forman curvas cerradas.
\end{itemize}

Para una carga positiva,

\[
\boxed{\text{campo radial hacia afuera}}.
\]

Para una carga negativa,

\[
\boxed{\text{campo radial hacia adentro}}.
\]


\subsection{Conversión de cargas}

\[
\boxed{\SI{1}{\milli\coulomb}=\SI{1e-3}{\coulomb}}
\]

\[
\boxed{\SI{1}{\micro\coulomb}=\SI{1e-6}{\coulomb}}
\]

\[
\boxed{\SI{1}{\nano\coulomb}=\SI{1e-9}{\coulomb}}
\]

\[
\boxed{\SI{1}{\pico\coulomb}=\SI{1e-12}{\coulomb}}.
\]


\subsection{Procedimiento para varias cargas puntuales}

\begin{enumerate}
    \item Ubicar las cargas y el punto donde se desea calcular el campo.
    \item Determinar la distancia desde cada carga hasta el punto.
    \item Identificar el sentido de cada campo:
    \[
    q_i>0\Rightarrow\text{se aleja de }q_i,
    \qquad
    q_i<0\Rightarrow\text{apunta hacia }q_i.
    \]
    \item Calcular
    \[
    E_i=k\frac{|q_i|}{r_i^2}.
    \]
    \item Obtener las componentes
    \[
    E_{i,x}=E_i\cos\theta_i,
    \qquad
    E_{i,y}=E_i\sin\theta_i.
    \]
    \item Asignar los signos según el sentido de cada componente.
    \item Sumar:
    \[
    E_x=\sum_iE_{i,x},
    \qquad
    E_y=\sum_iE_{i,y}.
    \]
    \item Calcular
    \[
    E=\sqrt{E_x^2+E_y^2}.
    \]
    \item Determinar la dirección analizando el cuadrante.
\end{enumerate}


\subsection{Resumen de fórmulas}

\begin{align*}
E &= \frac{F}{q_0},\\
F &= |q|E,\\
F_x &= qE_x, \qquad F_y=qE_y,\\
E &= k\frac{|Q|}{r^2},\\
r_i &= \sqrt{(x-x_i)^2+(y-y_i)^2},\\
E_{i,x} &= E_i\cos\theta_i, \qquad E_{i,y}=E_i\sin\theta_i,\\
E_x &= \sum_iE_{i,x}, \qquad E_y=\sum_iE_{i,y},\\
E &= \sqrt{E_x^2+E_y^2},\\
\theta_r &= \tan^{-1}\left(\frac{|E_y|}{|E_x|}\right),\\
\lambda &= \frac{dq}{d\ell}, \qquad dq=\lambda\,d\ell,\\
\sigma &= \frac{dq}{dA}, \qquad dq=\sigma\,dA,\\
\rho &= \frac{dq}{dV}, \qquad dq=\rho\,dV,\\
E_x &= \int dE_x, \qquad E_y=\int dE_y,\\
E_{\text{anillo},x} &= k\frac{Qx}{(x^2+a^2)^{3/2}}.
\end{align*}